\chapter{Heaps, priority queues}\label{chap:heaps}

Heaps are tree-based data structures that satisfy the heap property: parent nodes compare less than (or greater than) their children. They are efficient for implementing priority queues.

\section{Design rationale}

We implemented a binary heap using a dynamic array. This implicit representation (where children of the node at index $i$ are at $2i+1$ and $2i+2$) is memory efficient and provides excellent cache locality.

\section{Binary heap interface}

\begin{code}
  \lstinputlisting[style=C,linerange={3-12}]{libellul/include/libellul/type/heap/interface.h}
  \caption{The generic heap interface.}
  \label{code:heap:interface}
\end{code}

\section{Priority queue interface}

The priority queue is implemented as a wrapper around the binary heap. It allows users to manage a collection of elements where the "highest priority" element is always served first.

\section{Example variant: Rank-pairing heaps}

While binary heaps are excellent for general use, other variants like rank-pairing heaps offer better theoretical bounds for certain operations (like decrease-key or merge). However, due to the complexity and overhead of pointer-based structures, the binary heap often performs better in practice for simple priority queue needs.
